\section{Head-to-head: HCPC-v2 vs Self-RAG and CRAG}
\label{sec:headtohead}

The previous draft of this paper avoided direct comparison against
external retrieval-augmentation methods on the grounds that our
contribution is the failure-mode analysis rather than a new
state-of-the-art system. Reviewers correctly noted that this leaves the
question of whether HCPC-v2's advantage survives against published
selective-refinement systems unanswered. This section addresses that gap
by running Self-RAG and CRAG on the same query set, with the same
retrieval corpus and top-$k$, alongside our four conditions.

\subsection{Conditions}

Five conditions are evaluated end-to-end on the multidataset query set:

\begin{description}
  \item[\textsc{Baseline}] fixed 1024-token chunks, standard top-$k$ = 3
        retrieval, no refinement, Mistral-7B (Ollama) generator.
  \item[\textsc{HCPC-v1}] OR-gate selective refinement on the same
        retrieval. Always-refine when either similarity or cross-encoder
        gate triggers.
  \item[\textsc{HCPC-v2}] AND-gate + rank protection + merge-back, as
        described in §\ref{sec:method}.
  \item[\textsc{CRAG}] our reimplementation of \citet{yan2024crag}. The
        published evaluator is a private fine-tuned T5-large; we
        substitute a cross-encoder
        (\texttt{ms-marco-MiniLM-L-6-v2}) and threshold its logits to
        recover the three-bucket confidence signal (\texttt{correct} /
        \texttt{ambiguous} / \texttt{incorrect}). Knowledge refinement
        decomposes each retained passage into sentence-level strips and
        keeps the top-$k$ scoring strips. Web-search augmentation is
        disabled by default (closed-corpus comparison); the substitution
        and the disabled web path are documented transparently in
        \texttt{src/crag\_retriever.py}.
  \item[\textsc{Self-RAG}] the published
        \texttt{selfrag/selfrag\_llama2\_7b}
        checkpoint~\citep{asai2024selfrag}, fed the same retrieval
        passages and prompted with the model's training-time format
        (\texttt{[Retrieve]<paragraph>...</paragraph>}). Reflection
        tokens (\texttt{[Relevant]} / \texttt{[Fully supported]} /
        \texttt{[Utility:N]}) are parsed and reported as diagnostics.
        Self-RAG inference requires a CUDA GPU; the harness runs on M4
        for the four other conditions and is re-invoked on a GPU host to
        fill in the Self-RAG rows.
\end{description}

All conditions share: the same six datasets (§\ref{sec:multidataset}),
the same retrieval corpus, the same NLI-based faithfulness scorer, and
identical top-$k$ = 3.

\subsection{Preregistered reporting plan}
\label{sec:headtohead_plan}

For each of the six datasets we will report two paired quantities per
condition: the NLI-measured faithfulness rate and the binary
hallucination rate. That is a $6 \times 2 \times 5$ grid of cells,
populated from \texttt{results/headtohead/summary.csv}, with the
best condition per (dataset, metric) marked in bold and 95\% bootstrap
confidence intervals reported alongside the point estimate. All five
conditions share the same retrieval corpus, top-$k$, and NLI scorer so
that differences are attributable to the refinement policy rather than
the upstream pipeline. We commit to releasing per-query records so that
any conditional metric a reader wants (e.g., faithfulness on the subset
where Self-RAG emits \texttt{[Fully supported]}) can be recomputed
without re-running inference.

\subsection{Diagnostics}

Beyond aggregate metrics, the harness logs condition-specific decision
traces to \texttt{results/headtohead/diagnostics.jsonl}:

\begin{itemize}
  \item \textbf{CRAG:} per-passage labels and the resulting query-level
        decision (\texttt{correct} / \texttt{ambiguous} /
        \texttt{incorrect}), the number of strips retained, and the
        evaluator logits. These let us inspect whether CRAG's evaluator
        triggers on the same queries our coherence signals flag.
  \item \textbf{Self-RAG:} reflection-token counts per query and the
        parsed \texttt{[Relevant]} / \texttt{[Fully supported]} /
        \texttt{[Utility:N]} judgements. We report the conditional
        hallucination rate when Self-RAG self-judges its output as
        \texttt{Fully supported} versus when it does not — a measure of
        the calibration of the model's own faithfulness self-assessment.
\end{itemize}

\subsection{What HCPC-v2 has to beat}

We will treat the comparison as a success on three explicit conditions:

\begin{enumerate}
  \item HCPC-v2's faithfulness is within 1~pp of Self-RAG's on the
        domain-matched datasets (SQuAD, NaturalQs, TriviaQA), \emph{and}
        clearly better than CRAG-with-disabled-web.
  \item HCPC-v2's coherence-paradox signal (faith\_drop on HCPC-v1
        relative to baseline) reproduces the same direction we report on
        the original 2-dataset study, on at least 4 of 6 datasets.
  \item HCPC-v2's Self-RAG-style \emph{self-support rate} (the fraction
        of queries on which Self-RAG itself emits \texttt{[Fully
        supported]}) is correlated positively with our CCS metric on
        matched queries, validating that CCS captures something the
        Self-RAG model also tracks internally.
\end{enumerate}

\paragraph{Honest negative outcomes.} If Self-RAG dominates HCPC-v2 on
all six datasets, the paper's contribution narrows to the failure-mode
diagnosis (which is independently valuable but does not justify HCPC-v2
as a deployment recommendation); we will state this outcome plainly in
the discussion. If CRAG with web search disabled is competitive with
Self-RAG, the practical implication is that retrieval refinement matters
more than reflection-token training, which would also be worth reporting
as written.
